% Résumé
\begin{sommaire}{Langue des Signes Malgache, reconnaissance de signes, MediaPipe, GRU, application web}
Ce mémoire présente la reconnaissance de dix signes isolés de la Langue des Signes Malgache (LSM) et leur restitution en texte et en voix malgache. Quatre signeurs ont enregistré 400 vidéos ; les séquences d’un d’eux sont réservées au test. Chaque vidéo est convertie en points de squelette (pose et mains) extraits par MediaPipe, puis normalisée. Deux réseaux récurrents, le GRU et le LSTM, ont été entraînés et comparés à hyperparamètres identiques. Le LSTM atteint 82,0 \% d’exactitude sur le signeur de test et le GRU expérimental 92,0 \%. Le modèle final GRU, réentraîné sous la contrainte de conversion TFLite, atteint 94,0 \% ; conservant les points du corps et des mains, il a été converti au format léger TFLite (artefact de 247 Ko) et intégré à une application web : un backend FastAPI assure la reconnaissance à partir d’un court enregistrement webcam, tandis que le navigateur génère la prononciation malgache du signe reconnu. Dix essais vidéo sur l’application confirment le fonctionnement de la chaîne complète. Le prototype couvre un vocabulaire limité de signes isolés ; l’élargissement du corpus, puis la traduction de signes continus, constituent les prochaines étapes. Les résultats reposent sur un seul signeur de test et ne préjugent pas de performances générales.

\end{sommaire}

% Abstract
\begin{abstract}{Automatic recognition of isolated Malagasy Sign Language signs}{Malagasy Sign Language, sign recognition, MediaPipe, GRU, web application}
This thesis presents the recognition of ten isolated Malagasy Sign Language (LSM) signs and their restitution as text and spoken Malagasy. Four signers recorded 400 videos, and one signer is held out for testing. Each video is converted into pose and hand skeleton points extracted with MediaPipe and then normalised. Two recurrent networks, GRU and LSTM, were trained and compared under identical hyperparameters: the LSTM reaches 82.0\% accuracy on the test signer and the experimental GRU 92.0\%. The final GRU model, retrained under the TFLite conversion constraint, reaches 94.0\% accuracy; keeping both body and hand landmarks, it was converted to the lightweight TFLite format (a 247 KB artifact) and is deployed in a web application: a FastAPI backend performs the recognition from a short webcam recording, while the browser generates the Malagasy pronunciation of the recognized sign. Ten video tests on the application confirm that the complete chain works. The prototype covers a limited vocabulary of isolated signs; widening the corpus, and then translating continuous signs, are the next steps. The results rely on a single test signer and do not reflect general performance.

\end{abstract}
